% مطالبی از امید ریاضی که برای جلسه‌ای بعد کنار گذاشته شده‌اند (در هیچ jalase فعلاً \input نشده‌اند).

%% ------------------------------------------------------------
\section{چند مبحث دیگر از امید ریاضی}
\label{sec:expect-more}
%% ------------------------------------------------------------

\subsection{فرمول \lr{tail} برای امید ریاضی}

\begin{proposition}[label=prop:tail-sum]{امید ریاضی از روی \lr{tail}}
(الف) اگر $X$ فقط مقادیر $0,1,2,\dots$ را بگیرد: \ $\displaystyle\mathbb E[X]=\sum_{k=0}^{\infty}\P(X>k)$.\\
(ب) اگر $X\ge0$ پیوسته باشد: \ $\displaystyle\mathbb E[X]=\int_0^\infty\P(X>x)\,dx$.
\end{proposition}

\begin{proof}
(الف) هر عدد صحیح نامنفی $m$ برابر تعداد $k$هایی از $\{0,1,2,\dots\}$ است که $k<m$؛ پس
\[ X=\sum_{k=0}^{\infty}\ind\{X>k\}. \]
با خطی بودن (که برای جمع نامتناهیِ جمله‌های نامنفی هم برقرار است) و $\mathbb E[\ind_A]=\P(A)$ نتیجه می‌شود.\\
(ب) نسخهٔ پیوستهٔ همین ایده است:
\[ X=\int_0^\infty\ind\{X>x\}\,dx. \]
\end{proof}

این فرمول اغلب از تعریف ساده‌تر است، چون \lr{tail} بسیاری از توزیع‌ها فرم ساده‌ای دارد:
\begin{itemize}
  \item \textbf{هندسی:} $\P(X>k)=(1-p)^k$؛ پس $\mathbb E[X]=\sum_{k\ge0}(1-p)^k=\frac1p$، بدون مشتق‌گیری از سری.
  \item \textbf{نمایی:} $\P(X>x)=e^{-\lambda x}$؛ پس $\mathbb E[X]=\int_0^\infty e^{-\lambda x}dx=\frac1\lambda$، بدون مشتق‌گیری.
\end{itemize}

\subsection{امید ریاضی شرطی و قضیهٔ امید کل}

همان‌طور که احتمال شرطی داشتیم (جلسهٔ ۴)، می‌توانیم میانگینِ $X$ را \emph{به شرط} رخ دادن یک پیشامد $A$ (با $\P(A)>0$) حساب کنیم: کافی است در تعریف امید ریاضی، توزیع شرطی را بگذاریم.

\begin{definition}[label=def:cond-expect]{امید ریاضی شرطی به شرط یک پیشامد}
\[ \mathbb E[X\mid A]=\sum_xx\,\P(X=x\mid A)\qquad\text{\rl{یا}}\qquad \mathbb E[X\mid A]=\int_{-\infty}^{\infty}x\,f_X(x\mid A)\,dx. \]
\end{definition}

\begin{theorem}[label=thm:total-expect]{قضیهٔ امید کل}
اگر $A_1,\dots,A_n$ یک افراز از $\Omega$ باشند و $\P(A_i)>0$:
\[ \mathbb E[X]=\sum_{i=1}^n\P(A_i)\,\mathbb E[X\mid A_i]. \]
\end{theorem}

\begin{proof}
(حالت گسسته) طبق قانون احتمال کل، $p_X(x)=\sum_i\P(A_i)\,\P(X=x\mid A_i)$. دو طرف را در $x$ ضرب و روی $x$ جمع کنید و ترتیب دو جمع را عوض کنید.
\end{proof}

این همان قانون احتمال کل است، این بار برای میانگین: \key{میانگینِ کل، میانگینِ وزن‌دارِ میانگین‌های حالت‌های مختلف است}.

\begin{example}[label=ex:commute]{رفتن به دانشگاه}
دانشجویی با احتمال $0.3$ با اتوبوس به دانشگاه می‌رود که به طور متوسط ۲۰ دقیقه طول می‌کشد، و با احتمال $0.7$ پیاده می‌رود که به طور متوسط ۴۰ دقیقه طول می‌کشد. زمان متوسط رسیدن او چقدر است؟
\begin{solution}
$\mathbb E[T]=0.3\times20+0.7\times40=6+28=34$ دقیقه.
\end{solution}
\end{example}

\begin{example}[label=ex:geom-condition]{میانگین هندسی با شرط گذاشتن روی پرتاب اول}
اگر $X\sim\mathrm{Geom}(p)$، با شرط گذاشتن روی نتیجهٔ آزمایش اول، $\mathbb E[X]$ را بدون هیچ سری‌ای حساب کنید.
\begin{solution}
اگر آزمایش اول موفق باشد ($\P=p$)، $X=1$. اگر شکست بخورد ($\P=1-p$)، یک آزمایش «هدر» رفته و طبق بی‌حافظگی (جلسهٔ ۶) از نو شروع می‌کنیم؛ پس $\mathbb E[X\mid\fa{شکست اول}]=1+\mathbb E[X]$. طبق قضیهٔ امید کل:
\[ \mathbb E[X]=p\cdot1+(1-p)\bigl(1+\mathbb E[X]\bigr)\quad\Longrightarrow\quad p\,\mathbb E[X]=1\quad\Longrightarrow\quad \mathbb E[X]=\frac1p. \]
\end{solution}
\end{example}

\begin{example}[label=ex:exp-condition]{بی‌حافظگی و میانگین}
اگر $X\sim\mathrm{Exp}(\lambda)$ و $s>0$، امید ریاضی $\mathbb E[X\mid X>s]$ را بیابید.
\begin{solution}
در جلسهٔ ۸ دیدیم چگالی شرطی $f_X(x\mid X\ge s)=\lambda e^{-\lambda(x-s)}$ برای $x\ge s$ است؛ یعنی همان چگالی نمایی که به اندازهٔ $s$ به راست رفته. پس با تغییر متغیر $t=x-s$:
\[ \mathbb E[X\mid X>s]=\int_s^\infty x\,\lambda e^{-\lambda(x-s)}\,dx=\int_0^\infty(s+t)\,\lambda e^{-\lambda t}\,dt=s+\frac1\lambda. \]
اگر تا لحظهٔ $s$ صبر کرده باشیم، به طور متوسط \emph{باز هم} $\frac1\lambda$ دیگر باید صبر کنیم.
\end{solution}
\end{example}

\subsection{نامساوی مارکوف}

اگر فقط میانگین یک متغیر \emph{نامنفی} را بدانیم، باز هم می‌توانیم احتمال بزرگ شدنش را محدود کنیم:

\begin{theorem}[label=thm:markov]{نامساوی مارکوف}
اگر $X\ge0$ و $a>0$:
\[ \P(X\ge a)\le\frac{\mathbb E[X]}{a}. \]
\end{theorem}

\begin{proof}
برای هر $\omega$: اگر $X\ge a$ آنگاه $a\,\ind\{X\ge a\}=a\le X$، و وگرنه $a\,\ind\{X\ge a\}=0\le X$. پس $a\,\ind\{X\ge a\}\le X$. با گرفتن امید ریاضی از دو طرف (یکنوایی) و $\mathbb E[\ind\{X\ge a\}]=\P(X\ge a)$:
$a\,\P(X\ge a)\le\mathbb E[X]$.
\end{proof}

\begin{example}[label=ex:markov]{کران از روی میانگین}
(الف) میانگین نمرهٔ یک کلاس (از ۲۰) برابر ۱۲ است. حداکثر چه کسری از کلاس می‌توانند نمرهٔ ۲۰ گرفته باشند؟ (ب) یک سرور به طور متوسط ۴ درخواست در ثانیه دریافت می‌کند. احتمال اینکه در یک ثانیه دست‌کم ۲۰ درخواست برسد حداکثر چقدر است؟
\begin{solution}
(الف) $\P(X\ge20)\le\frac{12}{20}=0.6$. (ب) $\P(X\ge20)\le\frac4{20}=0.2$. اگر بدانیم تعداد درخواست‌ها $\mathrm{Poisson}(4)$ است، احتمال واقعی حدود $10^{-8}$ است؛ کران مارکوف ضعیف است چون جز میانگین هیچ اطلاعاتی از توزیع استفاده نمی‌کند. با دانستن واریانس (جلسهٔ بعد) کران‌های بهتری به دست می‌آید.
\end{solution}
\end{example}

\begin{moreinfo}{پارادوکس سن‌پترزبورگ}
سکهٔ سالمی را آن‌قدر می‌اندازیم تا شیر بیاید؛ اگر این کار $X$ پرتاب طول بکشد، $2^X$ تومان جایزه می‌گیرید. حاضرید برای شرکت در این بازی چقدر بپردازید؟ امید ریاضی جایزه
\[ \mathbb E[2^X]=\sum_{k=1}^{\infty}2^k\cdot\frac1{2^k}=1+1+1+\cdots=\infty \]
است؛ پس طبق «میانگین»، باید حاضر باشید هر مبلغی بپردازید! ولی تقریباً هیچ‌کس حاضر نیست بیش از چند هزار تومان بدهد، چون با احتمال $\tfrac78$ جایزه حداکثر ۸ تومان است. این پارادوکس (دانیل برنولی، ۱۷۳۸) نشان می‌دهد که وقتی امید ریاضی بی‌نهایت است یا \lr{tail} سنگین است، «میانگین» تصویر خوبی از رفتار متغیر نمی‌دهد.
\end{moreinfo}

%% تمرین‌ها

\begin{exercise}
در یک فروشگاه، ۶۰٪ مشتری‌ها به طور متوسط ۲ کالا و ۴۰٪ به طور متوسط ۷ کالا می‌خرند. هر مشتری به طور متوسط چند کالا می‌خرد؟
\answer{$0.6\times2+0.4\times7=4$}
\end{exercise}

\begin{exercise}
زمان پاسخ یک سرور (به میلی‌ثانیه) متغیری نامنفی با میانگین ۵۰ است. با نامساوی مارکوف، احتمال اینکه زمان پاسخ دست‌کم ۵۰۰ میلی‌ثانیه باشد را از بالا محدود کنید.
\answer{$\le0.1$}
\end{exercise}
